\capitulo{2}{Objetivos del proyecto}

A continuación, se detallan los objetivos que han guiado la realización del proyecto, diferenciando entre el objetivo general, los objetivos propios del \textit{software}, los objetivos técnicos y los objetivos personales.

\section{Objetivo general}

El objetivo general del proyecto es desarrollar una plataforma \textit{web} modular y extensible capaz de ejecutar flujos de procesamiento de imágenes mediante inteligencia artificial, devolviendo al usuario resultados visuales y datos estructurados que faciliten la interpretación del análisis realizado.

\section{Objetivos del \textit{software}}

\begin{itemize}
	\item Permitir que los usuarios se registren, inicien sesión y accedan a un panel propio.
	\item Ofrecer una interfaz desde la que el usuario pueda consultar los \textit{pipelines} disponibles y seleccionar el análisis que desea realizar.
	\item Permitir la subida de imágenes y la ejecución del procesamiento correspondiente en el servidor.
	\item Mostrar al usuario los resultados obtenidos de forma clara, incluyendo la imagen procesada y la información generada por el análisis.
	\item Incorporar funciones de administración para gestionar usuarios, permisos, servicios disponibles y \textit{pipelines}.
\end{itemize}

\section{Objetivos técnicos}

\begin{itemize}
	\item Diseñar una arquitectura modular que separe la interfaz, la lógica de control, la lógica de procesamiento y la persistencia de datos.
	\item Definir un modelo de datos que permita representar usuarios, permisos, servicios contratados, modelos, modos de ejecución, \textit{pipelines} y resultados.
	\item Organizar la ejecución de los análisis de forma desacoplada, permitiendo añadir nuevos modelos y modos sin modificar el flujo principal del sistema.
	\item Gestionar de forma segura las imágenes recibidas y los resultados generados, evitando la exposición directa de rutas internas del servidor.
	\item Validar el funcionamiento del sistema mediante pruebas automatizadas sobre los componentes principales de la aplicación.
\end{itemize}

\section{Objetivos personales}

\begin{itemize}
	\item Aplicar de forma integrada los conocimientos adquiridos durante el grado en desarrollo \textit{web}, bases de datos, ingeniería del \textit{software} e inteligencia artificial.
	\item Profundizar en el diseño de aplicaciones modulares, mantenibles y preparadas para crecer de forma progresiva.
	\item Mejorar la capacidad de planificación, organización y documentación de un proyecto completo de desarrollo de \textit{software}.
	\item Reforzar el uso de pruebas como mecanismo para comprobar el correcto funcionamiento del sistema.
	\item Acercar el desarrollo del proyecto a una estructura propia de un entorno profesional.
\end{itemize}